\section{Including and Inputing other Files}
\vspace{10.5cm}
	\begin{docCommand}{input}{\brackets{\sl{file-name}}}
		Imports the commands from filename.tex into the target file; it's equivalent to typing all the commands from filename.tex right into the current file where the \cs{input} line is. 
	\end{docCommand}
	\begin{docCommand}{include}{\brackets{\sl{file-name}}}
		The \cs{include} macro is bigger and is supposed to be used with bigger amounts of content, like chapters, which people might like to compile on their own during the editing process.
	\end{docCommand}
If you would like to get a better understanding on the difference between the \cs{include} vs \cs{input} command, then I would suggest checking out the following website:\\ 
\url{https://tex.stackexchange.com/questions/246/when-should-i-use-input-vs-include}\\
\newpage
